\section{Validation sur essais élémentaires}
\label{sec-validation}

Le chapitre~3 de la thèse de Guo valide Solidity sur trois essais élémentaires : flexion trois points, essai brésilien, compression polyaxiale. Lisjak valide Y-Geo sur l'amorçage d'une fissure en aile et sur des essais de compression simple et brésiliens. Ce chapitre présente les bancs équivalents de rockim. Aucun n'a encore tourné en campagne : ils sont préparés, avec leurs critères écrits avant calcul, et n'ont tourné qu'en essai de fumée. Il présente ensuite les comparaisons plus anciennes à des simulations publiées et l'état de la calibration sur le Red Bohus.

\subsection{Bancs préparés}

\Cref{tab-bancs-prepares} résume les bancs préparés. Les références sont les mêmes que celles des codes de référence, avec un critère chiffré en plus.

\begin{table}[htbp]
\centering
\caption{Bancs de validation élémentaire préparés, avec leur référence et leur critère principal.}
\label{tab-bancs-prepares}
\small
\begin{tabularx}{\linewidth}{@{}l X X X@{}}
\toprule
banc & référence & critère principal & essai de fumée \\
\midrule
fissure pressurisée (P3.1) & Guo §{2.4} ; mécanique de la rupture ; fissure cohésive exacte & convergence à $5~\%$ ; écart à Guo $\leq 15~\%$ & amorçage $1{,}91$-$1{,}99$~MPa \\
fissure en aile (P3.3) & Lisjak : $\sigma_{1c} = 7{,}0$~MPa, $\theta = 64^\circ$ ; mécanique de la rupture 5,1 à 14,6~MPa & $|\sigma_{1c}/7{,}0 - 1| \leq 15~\%$, angle à $10^\circ$ & chaîne fonctionnelle, valeurs sans objet \\
flexion trois points (P4.1) & Guo §{3.3} ; poutre de Timoshenko & raideur et contrainte de rupture & chaîne fonctionnelle \\
brésilien (P4.2) & Guo §{3.4} ; Hondros & $f_{bt}/f_t$ entre $0{,}85$ et $1{,}15$ & chaîne fonctionnelle \\
élasticité (P1.2, P1.3) & Kirsch ; champ homogène & erreur nodale $\leq 10^{-6}$ ; $K_t$ à $2~\%$ & champ homogène à $7{,}5\times10^{-8}$ \\
impact de bouton (B8) & Saksala 2011 & force maximale à $15~\%$ & $-14{,}7~\%$ sur maillage grossier \\
\bottomrule
\end{tabularx}
\end{table}

Deux bancs méritent une remarque. Pour le brésilien, les critères écrits avant calcul annoncent qu'ils ne peuvent pas passer ensemble : Guo trouve $f_{bt} = 1{,}96$~MPa pour $f_t = 3$~MPa, soit $0{,}65\,f_t$, qu'il attribue à la rupture en cisaillement sous les plateaux plans, alors qu'un modèle qui respecte Hondros doit amorcer au centre près de $f_t$. Pour la fissure en aile, deux des quatre solutions de mécanique de la rupture de la fourchette de Lisjak sont des reconstructions non vérifiées. La compression polyaxiale de Guo n'a pas de banc dans la campagne ; Guo y juge l'orientation des plans de cisaillement à l'œil, contre l'angle $\pm(45^\circ + \varphi/2)$.

\subsection{Comparaisons antérieures à des simulations publiées}

Trois comparaisons sont antérieures à la campagne. Elles portent toutes sur la simulation calée d'un autre code et non sur un essai indépendant.

Sur la compression simple 2D avec insertion adaptative de Yan et al.~\cite{yan2023adaptive}, rockim retrouve la résistance de 51~MPa. Relancée avec le binaire courant, la même configuration donne $51{,}43$~MPa contre $51{,}07$~MPa auparavant, soit $+0{,}70~\%$ pour une tolérance de $0{,}29~\%$, et $34~\%$ de joints rompus en plus ; la cause, évolution du code ou maillage de Voronoï dépendant de la bibliothèque standard, n'est pas départagée.

Sur la fracturation hydraulique en paroi de forage d'AbuAisha et al.~\cite{abuaisha2017hm}, calculée en 2D avec un couplage hydromécanique sans écoulement, la pression de rupture est trop forte de $13{,}2~\%$ en isotrope et de $25~\%$ en anisotrope ; l'ouverture d'une fissure développée est à $-4~\%$ de Sneddon et le déplacement de paroi à $-2{,}0~\%$ de Lamé.

Une comparaison à Abaqus avec la loi de Saksala sur le même maillage, sphère à 8~m/s, donnait le 31 juillet 2026 un pic de force filtré à $+7{,}9~\%$ et une énergie à $+0{,}9~\%$, avec des critères fixés avant calcul. Les fichiers de cette comparaison ne sont pas dans le dépôt.

\subsection{Calibration sur le granite de Red Bohus}

Les codes de référence calent leurs paramètres sur des essais de compression simple et brésiliens, en quatre familles (numériques, volume, interaction élastique, rupture), selon la procédure de Tatone et Grasselli~\cite{tatone2015}. La calibration de rockim sur le Red Bohus (résistance en compression $126{,}6$~MPa, brésilienne $10{,}27$~MPa, module 77,7~GPa~\cite{dumoulin2024data}) n'est pas conclue. Un premier jeu bayésien 2D à grains reproduisait la traction à $-6{,}7~\%$ et la compression à $+9{,}9~\%$ ; il a été invalidé par la correction des amortissements. Un second jeu, par émulateur gaussien, se trompe sur le triaxial hors calage de $+36$ à $+53~\%$ entre 50 et 100~MPa de confinement, l'enveloppe de Mohr-Coulomb linéaire des joints surestimant les forts confinements. La partie de la campagne prévue pour cette validation (calage sur compression, brésilien et triaxial à 20 et 50~MPa, puis prédiction sans retouche à 75 et 100~MPa) n'est pas ouverte. Aucune calibration du Red Bohus n'alimente aujourd'hui la FDEM 3D.
